\section{Secure Shared Memory}
\label{sec:sharedmemory}

For an enclave, communicating with the outside world is an essential capability.
Indeed, very few applications can be reduced to a fully self-contained binary that does not require external input, interactions with the outside, or the capability to output more than a few bytes of data.
Developing efficient and secure enclave communication mechanisms is essential.

\subsection{Transient Attack on Shared Memory}

With transient execution attacks, a program that manages sensitive data can be coerced into leaking private data speculatively through a microarchitectural side channel (see Section \ref{sec:specandside}).
This means that secrets can potentially be extracted from an enclave and leaked to an adversary like a malicious OS or a co-running program or enclave.
One way to defend against these types of attacks is to close the microarchitectural side channels.
This way, even if a secret is accessed speculatively, it cannot be exfiltrated outside of the victim's security domain.
This can be achieved through strong microarchitectural isolation between security domains (see Section \ref{sec:hardwaremechanisms}).

Unfortunately, in the case of shared memory, things become more complicated.
If two security domains share a region of physical memory, they irremediably entangle their microarchitectural state through shared caches and memory coherence mechanisms.
This shared state can then be used to mount a Spectre-style speculative attack.
We implemented such an attack on an insecure version of our platform where shared memory is enabled but when no defense mechanism is implemented.
Our insecure platform is vulnerable to such attacks, and the attack is mitigated by our defense mechanism.

\subsection{Tentative Solution: Secure Drawers}
\label{sec:secdrawer}
Our first attempt at a new communication primitive is called secure drawer.
Very much like the ones used at bank counters, a secure drawer is a memory region that can be filled up by one security domain and then pushed and pulled by another security domain to access it.
This primitive has clear security boundaries as no security domains have concurrent access to the same memory region.
Nevertheless, it has important performance drawbacks as it can only manipulate coarse memory regions as defined by the underlying physical memory partitioning primitive (32MB in our case) and several SM calls are required to push and pull the memory region.
It also requires the SM to flush the LLC partition when moving the drawer and to copy all the data used by the enclave to the secure drawer.
Finally, the enclave programmer needs to modify their application to deal with the new primitive and potentially different address spaces, move all the require data around, and handle concurrency correctly.
In the end, secure drawers are far from transparent to the programmer and impose important performance limitations as we show in Section \ref{sec:eval}.


\subsection{New Defense Strategy}

Because accesses to shared memory will always be visible from one security domain to another, we decided to block the enclave from speculatively accessing shared memory.
This prevents the enclave from transmitting the secret through any side channel that exploits a shared microarchitectural structure in the memory hierarchy which is the only shared microarchitecture between two cores running concurrently.
Any access to shared memory will have to wait and become non-speculative before being dispatched to the memory execution pipeline.
This guarantees that any access to shared memory is performed respecting the enclave's program semantics and that no speculative execution can leave any trace in the microarchitecture of shared memory hierarchy.

Note that our mechanism only protects \emph{enclaves} from other malicious programs running on the same machine.
Without modifications, the OS will not use our trusted memory mechanism and therefore will not benefit from our new defense.
As a result, it is still vulnerable to transient execution attacks.

\textbf{Non-cacheable Mechanism}
Note that modern processors already include a mechanism to prevent some form of speculative memory accesses, namely non-cacheable memory accesses.
For instance, MMIO accesses should never be performed speculatively, and so they are declared non-cacheable.
Unfortunately, these mechanisms %are not designed for security and 
cannot be repurposed for secure shared memory.
Indeed, the virtual address translation of non-cached accesses is still typically done speculatively (including in the prototype processor we used).
Speculative address translation accesses are known to be enough to leak secrets, so we need to design a mechanism fully preventing any form of shared-memory speculation.

\subsection{Dual Page Table Mechanism}

\label{sec:vmanddualpt}

\begin{figure}[t]
\centering
\includegraphics[width=\columnwidth]{figures/dual_page_table.png}
\caption{The Dual Page Table Mechanism}
%\vspace{8pt}
\label{fig:dual_pt}
\end{figure}

Sanctum \cite{costan2016sanctum} already presents hardware mechanisms for memory isolation in the context of an untrusted OS.
Like RISC-V Physical Memory Protection (PMP) \cite{waterman2015risc}, Sanctum introduces a mechanism to partition physical memory and prevent a malicious OS to access another security domain's private memory regions.
 At its core, the mechanism uses two 64-bit registers per core (one for trusted memory and one for untrusted memory) as memory bitmaps.
Each bit represents a same size physical memory region (typically 32MB) and functions as an access-rights flag for the currently-running program. Citadel borrows this mechanism for its simplicity.

Sanctum also uses a dual-page-table mechanism (cf. Figure \ref{fig:dual_pt}) to differentiate between trusted and untrusted memory based on virtual addresses, which we also borrow in Citadel. 
Two new registers \texttt{TMVBASE} and \texttt{TMVMASK} are used to define \texttt{TMVRANGE} which corresponds to the trusted memory virtual range (the enclave's private memory in general).
The SM enforces as an invariant that the physical memory mapped by \texttt{TMVRANGE} is owned by the corresponding security domain.
When a memory access to a virtual address is performed by a program, one of the two page tables needs to be picked to translate the virtual address.
If the virtual address falls in \texttt{TMVRANGE}, it corresponds to an address within trusted memory and the translation is performed using the \emph{trusted page table}.
The trusted memory page table root is stored in another new register called \texttt{TMPTBR} accessible to the page walker.
Another important invariant is that the trusted enclave page table is stored within trusted enclave memory.
This is at first enforced by the SM during setup and then by the enclave itself.
If the virtual address falls outside of \texttt{TMVRANGE}, it corresponds to an address in untrusted memory and the translation is performed using the \emph{untrusted page table}.
The untrusted page table is managed by the OS and stored in untrusted memory.
The root of the OS page table is stored in the standard \texttt{SATP} register.
Note that for every physical address used during address translation, additional bound checks are performed to ensure they belong to the same security domain as the virtual address.


The above mechanism is especially handy to enforce security in the presence of speculative side channels
as there is no need to perform any address translation to know if a memory access is trusted or not.

\subsection{Design Details}

In an out-of-order processor, the \textit{memory reservation station}s (or \MemRS{}) are structures used to queue and schedule different kinds of instructions, among which the memory requests. 
Following the register renaming stage, each memory instruction is simultaneously enqueued in both the reorder buffer and the appropriate reservation station. 
In the base design, the reservation station selects which entry to dispatch to the memory execution pipeline (\MemExePipeline{}) based on its operands' availability. 
During dispatch, the entry is also dequeued from the \MemRS{}, at which point the \MemExePipeline{} begins processing the request. 
The \MemExePipeline{} contains several stages 
including Dispatch, Register Read, Translation Execution, and Finish, after which the instruction is placed in the load/store queue to be scheduled and sent to memory.

Our approach is to identify accesses to untrusted memory early in the \MemExePipeline{} and delay the translation and processing of the memory accesses until the instruction is no longer speculative.
These changes require the modification of \MemRS{} and \MemExePipeline{}. 

\begin{figure}[t]
\centering
\includegraphics[width=\columnwidth]{figures/MemExePipeline.png}
\caption{Changes to Memory Execution Pipeline. New elements are shown in blue.}
%\vspace{8pt}
\label{fig:MemExe}
\end{figure}

\textbf{Modifying MemExe}
Logic to perform security checks and send back shared-memory speculative requests is added to \MemExePipeline{} (cf. Figure \ref{fig:MemExe}).
First, we add logic in the translation execution stage to determine if the request is accessing the enclave's trusted memory or untrusted shared memory using \texttt{TMVBASE} and \texttt{TMVMASK}. 
If the request is accessing untrusted memory, an additional check is performed to see whether the access is speculative or not.
The top entry from the ROB is compared to the tag of the request in the pipeline. 
If the tag matches, this indicates a \textit{non-speculative} access, since all other previous instructions in the program have been committed.
In that case, the request is dispatched to the next stage of the pipeline.
Whenever an access to shared memory is found to be speculative, the request is squashed, and a signal is sent to the \MemRS{} indicating that it should re-dispatch the request once it reaches the head of the ROB.
The request may never reach the head of the ROB if there is rollback due to mis-speculation.

\textbf{Decorrelating Request Dispatching and Request Dequeuing}
Our mechanism requires the \MemRS{} to retain the accesses to shared memory in order to re-dispatch them once they become non-speculative.
The original design performs a \MemRS{} data dispatch and dequeue simultaneously once the request enters the \MemExePipeline{}. 
These two operations now need to occur at different stages in the pipeline: data is first dispatched from the \MemRS{} and is dequeued in the translation execution stage if and only if the request is determined to be secure. 
We modify the interface between \MemRS{} and \MemExePipeline{} to provide the correct interfaces and logic.
This includes augmenting the request data with the corresponding \MemRS{} entry index.

\textbf{Modifying \MemRS{}}
The \MemRS{} logic is modified to implement re-dispatching of requests once they become non-speculative.
First, several new status bits are added for each \MemRS{} entry, specifically to track if a) a given request is still to be dispatched for the first time, and b) if a given request needs to be re-dispatched. 
The first bit (\textit{to\_dispatch}) is necessary since entries can now still reside in the \MemRS{} in the time between being dispatched and dequeued (or being signaled for redispatch), and should not be selected for dispatch again.
The second bit (\textit{to\_redispatch}) is set by the \MemExePipeline{} when it squashes a speculative insecure request. 
We additionally add logic to compare each \MemRS{} entry to the entry at the head of the ROB, and raise the entry's \textit{dispatch\_ready} signal on a match.
Finally, the dispatch logic is modified to use the new status bits. Rather than simply checking for the oldest valid entry whose arguments are ready, an additional condition is added to check if either \textit{to\_dispatch} is set for the entry or \textit{to\_redispatch} and \textit{redispatch\_ready} are set for the given entry.

\textbf{Disabling the Mechanism for the OS}
Blocking speculation on memory access can have significant performance impact.
We want to make sure that the OS is still allowed to speculate on its own memory accesses when not using the trusted memory mechanism.
As a result, we disable our mechanism when \texttt{TMVRANGE} is set to a specific empty range.
Note that the OS could always be modified to take advantage of our trusted memory mechanism and protect itself from speculative side channel attacks.

\subsection{DMA For Shared Memory}

\subsubsection{Specter Variant Using The Shared L1}
Enabling shared memory without further precautions means that the victim private L1 cache would now be shared between data from the enclave private memory and data from shared memory.
This creates new, subtle cache side channels rooted in cache coherency.
Indeed, an attacker can exploit delays due to cache coherency mechanisms to observe if an element of shared memory is present in the victim L1 cache or not.
For instance, if the attacker tries to write a location in memory for which the victim L1 holds a copy, the operation will take longer as it will need to wait for the victim L1 copy to be invalidated.
Similarly, if the attacker tries to read a memory location for which the L1 holds a dirty copy, the operation will take longer as it will need to wait for the dirty copy to be brought back to the L2.
Unfortunately, the state of the L1 cache not only depends on non-speculative shared memory accesses, but it is also the (speculative) private memory accesses.
In particular, a copy of shared memory data could be evicted by a speculative access to private enclave memory.
This breaks our security invariant that microarchitectural state that is modifiable by enclave code running speculatively should not be observable by an attacker.
This makes it possible for the attacker to mount a variant of prime-and-probe and perform Spectre-style transient execution attacks.

\textbf{Tentative Solution: Partitioning the L1 Cache}

Using way partitioning to allocate a portion of the L1 cache to shared memory accesses seems to be a solution.
Unfortunately, this simple solution hides a lot of complexity as it would not only require partitioning the L1 but also adding state to keep track of which access is private or not and duplicating most of the logic of the L1 to enforce timing independence between the operations of the different partitions (e.g., when accessing cache metadata concurrently).
This is not only complex to implement, but could also impact the latency of each memory access, which would dramatically affect performance.

\textbf{Direct Memory Access for Shared Memory}
We decided to treat shared memory accesses differently and to not cache them in the L1.
Instead, when an access to shared memory is performed, we send the request through the core DMA channel. 
This solution involves arbitrating the different DMA requests between cores.
Note that DMA partitioning between cores was actually ignored in MI6, representing a security flaw as private TLB accesses are already performed through DMA accesses.
We implement this solution in our prototype.